(1) 理想混合液中溶剂化学势为 $\mu_w(T,p,x_w)=\mu_w^*(T,p)+RT\ln x_w$，其中摩尔分数 $x_w=n_w/(n_w+n_s)<1$。混合项 $RT\ln x_w<0$，因此同压时溶液中溶剂化学势低于纯溶剂。

(2) 取溶液侧压强比纯溶剂侧高 $\Pi$，近似不可压溶剂摩尔体积 $\bar v_w$ 常数。化学势差平衡 $\bar v_w\Pi+RT\ln x_w=0$，即 $\Pi=-(RT/\bar v_w)\ln x_w$。当 $n_s\ll n_w$，$-\ln x_w\simeq n_s/n_w$ 且 $V\simeq n_w\bar v_w$，所以 $\boxed{\Pi=(n_s/V)RT=c_sRT}$。

(3) 理想二元混合自由能为 $\Delta G_{\rm mix}=RT[n_w\ln x_w+n_s\ln x_s]<0$，其中 $x_s=1-x_w$。因混合自由能的熵项正且理想混合焓为零，$\Delta S_{\rm mix}=-R[n_w\ln x_w+n_s\ln x_s]>0$。稀溶质展开中，驱动力仍是溶质组分占据体积产生的配置熵。

(4) 若溶质完全解离成多个独立离子，理想渗透浓度需计全部粒子数，形成van't Hoff系数；若溶质也能透膜，则两侧溶质可交换，平衡条件不同且不再维持原渗透压。高浓度非理想液体必须以活度 $a_w$ 替代摩尔分数 $x_w$。
